% !PRINTED 587
\weberpage{587}
It remains to be shown that no two of the numbers (13) are equal to
one another. If we had
\[
\boldsymbol{s}\left(\frac{2\varrho}{\mu}\right) = \boldsymbol{s}\left(\frac{2\varrho'}{\mu}\right),
\]
without $\varrho \equiv \varrho'$ (mod $\mu$) holding, then by (5)
\[
\frac{2(\varrho + \varrho')}{\mu} \equiv 1 \quad (\mathrm{mod}\ 2, \omega)
\]
would have to hold. Hence $a(1+\omega)$ would have to be divisible by 2,
and hence also, if $\omega$ is conjugate to $\omega'$, $a(1+\omega')$ divisible by 2,
and consequently
\[
a(\omega - \omega') = \sqrt{\Delta} \equiv 0 \quad (\mathrm{mod}\ 2).
\]
Thus $\Delta$ would be divisible by 4, and $\frac{1}{4}\Delta = \frac{1}{4}a^2(\omega - \omega')^2$
would likewise still be a discriminant, which contradicts the assumption
that $\Delta$ is a fundamental discriminant\footnote{$b$ would have to be even and $a-b+c$ divisible by 4, and
$\frac{1}{4}\Delta = \left(\frac{b}{2}-a\right)^2 - a(a-b+c)$ would be a discriminant. For an arbitrary
discriminant $D$ the same end could be attained by adding to the
numerator of the expressions (13) a power of 2 as a factor.}. Accordingly the
quantities (13), whose number amounts to $m-1$ (Bd. II, § 165), are all
different from one another, and there results
\weq{14}{A(x) = \overset{\varrho}{\prod}\left[x - \boldsymbol{s}\left(\frac{2\varrho}{\mu}\right)\right].}

By § 44, (20) the function $\boldsymbol{s}(u)$ satisfies the equation:
\[
\boldsymbol{s}\left(u + \frac{\omega}{2}\right) = \frac{1}{\boldsymbol{s}(u)},
\]
and, because of (8) and (4):
\[
\boldsymbol{s}\left[\mu\left(u + \frac{\omega}{2}\right)\right] = \frac{(-1)^{cx}}{\boldsymbol{s}(\mu u)},
\]
and accordingly there results from (10), if $u$ is replaced by $u + \dfrac{\omega}{2}$:
\weq{15}{\frac{(-1)^{cx}}{\boldsymbol{s}(\mu u)} = \frac{1}{\boldsymbol{s}(u)}\,\frac{A\left(\dfrac{1}{\boldsymbol{s}}\right)}{D\left(\dfrac{1}{\boldsymbol{s}}\right)} = \frac{(-1)^{cx} D(\boldsymbol{s})}{\boldsymbol{s}(u) A(\boldsymbol{s})},}
\weq{16}{A(\boldsymbol{s}) A\left(\frac{1}{\boldsymbol{s}}\right) = (-1)^{cx} D(\boldsymbol{s}) D\left(\frac{1}{\boldsymbol{s}}\right).}
% [source sic: equations (15)--(16) behave as though the external factor \mu in (10) has been absorbed into a normalized A]
